% Options for packages loaded elsewhere
\PassOptionsToPackage{unicode}{hyperref}
\PassOptionsToPackage{hyphens}{url}
\PassOptionsToPackage{dvipsnames,svgnames,x11names}{xcolor}
\documentclass[
  11pt,
]{article}
\usepackage{xcolor}
\usepackage[margin=2.5cm]{geometry}
\usepackage{amsmath,amssymb}
\setcounter{secnumdepth}{-\maxdimen} % remove section numbering
\usepackage{iftex}
\ifPDFTeX
  \usepackage[T1]{fontenc}
  \usepackage[utf8]{inputenc}
  \usepackage{textcomp} % provide euro and other symbols
\else % if luatex or xetex
  \usepackage{unicode-math} % this also loads fontspec
  \defaultfontfeatures{Scale=MatchLowercase}
  \defaultfontfeatures[\rmfamily]{Ligatures=TeX,Scale=1}
\fi
\usepackage{lmodern}
\ifPDFTeX\else
  % xetex/luatex font selection
    \setmainfont[]{Charter}
    \setmonofont[Scale=0.85]{Menlo}
  \setmathfont[]{STIX Two Math}
\fi
% Use upquote if available, for straight quotes in verbatim environments
\IfFileExists{upquote.sty}{\usepackage{upquote}}{}
\IfFileExists{microtype.sty}{% use microtype if available
  \usepackage[]{microtype}
  \UseMicrotypeSet[protrusion]{basicmath} % disable protrusion for tt fonts
}{}
\usepackage{setspace}
\makeatletter
\@ifundefined{KOMAClassName}{% if non-KOMA class
  \IfFileExists{parskip.sty}{%
    \usepackage{parskip}
  }{% else
    \setlength{\parindent}{0pt}
    \setlength{\parskip}{6pt plus 2pt minus 1pt}}
}{% if KOMA class
  \KOMAoptions{parskip=half}}
\makeatother
\usepackage{longtable,booktabs,array}
\usepackage{calc} % for calculating minipage widths
% Correct order of tables after \paragraph or \subparagraph
\usepackage{etoolbox}
\makeatletter
\patchcmd\longtable{\par}{\if@noskipsec\mbox{}\fi\par}{}{}
\makeatother
% Allow footnotes in longtable head/foot
\IfFileExists{footnotehyper.sty}{\usepackage{footnotehyper}}{\usepackage{footnote}}
\makesavenoteenv{longtable}
\ifLuaTeX
\usepackage[bidi=basic]{babel}
\else
\usepackage[bidi=default]{babel}
\fi
\babelprovide[main,import]{brazilian}
\ifPDFTeX
\else
\babelfont{rm}[]{Charter}
\fi
% get rid of language-specific shorthands (see #6817):
\let\LanguageShortHands\languageshorthands
\def\languageshorthands#1{}
\setlength{\emergencystretch}{3em} % prevent overfull lines
\providecommand{\tightlist}{%
  \setlength{\itemsep}{0pt}\setlength{\parskip}{0pt}}
\usepackage{bookmark}
\IfFileExists{xurl.sty}{\usepackage{xurl}}{} % add URL line breaks if available
\urlstyle{same}
\hypersetup{
  pdftitle={Plano do workshop},
  pdfauthor={Giovanni Vasconcelos · CESUPA Tech Summit 2026},
  pdflang={pt-BR},
  colorlinks=true,
  linkcolor={black},
  filecolor={Maroon},
  citecolor={Blue},
  urlcolor={blue},
  pdfcreator={LaTeX via pandoc}}

\title{Plano do workshop}
\usepackage{etoolbox}
\makeatletter
\providecommand{\subtitle}[1]{% add subtitle to \maketitle
  \apptocmd{\@title}{\par {\large #1 \par}}{}{}
}
\makeatother
\subtitle{Attention is All You Need: construindo um Transformer do zero}
\author{Giovanni Vasconcelos · CESUPA Tech Summit 2026}
\date{}

\begin{document}
\maketitle

\renewcommand*\contentsname{Sumário}
{
\hypersetup{linkcolor=}
\setcounter{tocdepth}{3}
\tableofcontents
}
\setstretch{1.15}
\section{1. Visão geral}\label{visuxe3o-geral}

\textbf{Objetivo.} Ao final das duas horas, cada participante terá
implementado os componentes centrais de um Transformer decoder-only
(atenção, multi-head attention e o bloco com conexões residuais) e usado
esse código para treinar um modelo de linguagem que gera partidas de
xadrez. O foco é a compreensão do mecanismo, não o desempenho do modelo.

\textbf{Público.} Estudantes de todos os semestres. O roteiro foi
desenhado para quem já teve contato com programação em Python e noções
de álgebra linear. Estudantes dos primeiros semestres acompanham a
intuição pelos slides e usam as células de recuperação do notebook para
não perder o fio. Espera-se uma turma de mais de 40 pessoas, com pelo
menos dois monitores.

\textbf{Formato.} Metade exposição, metade prática, sempre alternadas:
cada bloco teórico precede imediatamente o exercício correspondente.
Toda a parte prática acontece no Google Colab. O participante precisa
apenas de um navegador e de uma conta Google; nada é instalado nos
computadores do laboratório.

\textbf{Materiais.}

\begin{longtable}[]{@{}
  >{\raggedright\arraybackslash}p{(\linewidth - 2\tabcolsep) * \real{0.3889}}
  >{\raggedright\arraybackslash}p{(\linewidth - 2\tabcolsep) * \real{0.6111}}@{}}
\toprule\noalign{}
\begin{minipage}[b]{\linewidth}\raggedright
Arquivo
\end{minipage} & \begin{minipage}[b]{\linewidth}\raggedright
Função
\end{minipage} \\
\midrule\noalign{}
\endhead
\bottomrule\noalign{}
\endlastfoot
\texttt{workshop\_aluno.ipynb} & Notebook da turma, com seis exercícios
e células de verificação \\
\texttt{workshop\_gabarito.ipynb} & Mesmo notebook com as soluções; para
o ministrante e os monitores \\
\texttt{pretreino\_referencia.ipynb} & Treina o modelo de referência no
Colab (executar uma vez, antes do evento) \\
\texttt{mini\_gpt.py} & Implementação de referência, importada pelas
células de recuperação \\
\texttt{checagens.py} & Verificações automáticas dos exercícios \\
\texttt{tabuleiro.py} & Tabuleiro interativo (arrastar e soltar) para
jogar contra o modelo no Colab \\
\texttt{slides/attention\_workshop.pptx} & Apresentação (30 slides, com
notas do apresentador em todos) \\
\texttt{docs/GUIA\_DE\_ESTUDOS.md} & Guia de estudos para a preparação
do ministrante \\
\texttt{data/xadrez.txt.gz} & 60 mil partidas do Lichess (2013), já
filtradas e validadas \\
\end{longtable}

\section{2. Preparação}\label{preparauxe7uxe3o}

O workshop acontece em 14 de outubro de 2026 (D). As datas entre
parênteses são as correspondentes no calendário.

\subsection{D − 14 (30/09): repositório}\label{d-14-3009-reposituxf3rio}

\begin{enumerate}
\def\labelenumi{\arabic{enumi}.}
\item
  O repositório público já existe:
  \url{https://github.com/giovannibragasv/transformer-do-zero}. O
  usuário \texttt{giovannibragasv} já está configurado em
  \texttt{scripts/gerar\_notebooks.py}, \texttt{slides/gerar\_slides.js}
  e \texttt{README.md}.
\item
  Se algo mudar, regenere os artefatos:

\begin{verbatim}
.venv/bin/python scripts/gerar_notebooks.py
cd slides && npm install && node gerar_slides.js attention_workshop.pptx
\end{verbatim}
\item
  Fazer o commit de tudo, exceto \texttt{data\_raw/},
  \texttt{data/xadrez.txt} e \texttt{.venv/} (já listados no
  \texttt{.gitignore}).
\item
  Abrir
  \texttt{https://colab.research.google.com/github/giovannibragasv/transformer-do-zero/blob/main/workshop\_aluno.ipynb}
  e confirmar que o notebook abre e que a primeira célula clona o
  repositório.
\item
  Gerar um link curto (e, se possível, um QR code) para esse endereço e
  atualizar o slide 4.
\end{enumerate}

\subsection{D − 10 (04/10): modelo de
referência}\label{d-10-0410-modelo-de-referuxeancia}

\begin{enumerate}
\def\labelenumi{\arabic{enumi}.}
\tightlist
\item
  Abrir \texttt{pretreino\_referencia.ipynb} no Colab com GPU.
\item
  Executar as células. O treino dura 60 minutos e salva checkpoints a
  cada 10 minutos.
\item
  Baixar \texttt{xadrez\_referencia.pt} (cerca de 19 MB), colocar em
  \texttt{modelos/} e fazer o commit.
\item
  Anotar a taxa de legalidade impressa ao final. Ela entra na fala do
  slide 25 e serve de referência para a seção 3.4 do notebook.
\end{enumerate}

Se o treino for interrompido, o último checkpoint salvo já é utilizável.
Se o modelo não estiver disponível no dia, o notebook usa
automaticamente o modelo treinado pela turma.

\subsection{D − 7 (07/10): ensaio}\label{d-7-0710-ensaio}

\begin{enumerate}
\def\labelenumi{\arabic{enumi}.}
\tightlist
\item
  Com uma conta Google diferente da sua, sem nada salvo no Drive,
  executar o \texttt{workshop\_aluno.ipynb} do início ao fim como se
  fosse um participante. Resolver os exercícios sem consultar o gabarito
  e cronometrar cada parte.
\item
  Executar o \texttt{workshop\_gabarito.ipynb} inteiro e salvar a versão
  executada (com as saídas). Ela serve para demonstração caso o Colab
  fique instável no dia.
\item
  Ensaiar a exposição com os slides em voz alta. As notas do
  apresentador trazem o conteúdo de cada fala; o tempo de cada bloco
  teórico é de 10 minutos.
\item
  Fazer o estudo descrito no guia (\texttt{GUIA\_DE\_ESTUDOS.md}), em
  especial a seção de perguntas difíceis.
\end{enumerate}

\subsection{D − 3 (11/10): logística}\label{d-3-1110-loguxedstica}

\begin{itemize}
\tightlist
\item
  Confirmar com a organização o laboratório, o número de máquinas, o
  projetor e a conexão com a internet. Cada participante baixa cerca de
  30 MB na primeira célula; para 40 pessoas, são pouco mais de 1 GB em
  alguns minutos.
\item
  Confirmar que o navegador dos laboratórios permite login em contas
  Google.
\item
  Reunir os monitores por 30 minutos: percorrer o gabarito, combinar a
  divisão da sala e apresentar a lista de erros comuns (seção 5).
\item
  Preparar um cartaz ou slide com o link curto do notebook.
\end{itemize}

\subsection{Dia D (14/10)}\label{dia-d-1410}

\begin{itemize}
\tightlist
\item
  Chegar 30 minutos antes. Testar o projetor com os slides e com o Colab
  (fonte do notebook ampliada para 125\% ou mais).
\item
  Deixar abertos em abas: slides, notebook do aluno, gabarito executado
  e o notebook de gabarito pronto para a demonstração ao vivo do slide
  3.
\item
  Conectar-se a uma GPU no Colab antes do início, para a demonstração.
\end{itemize}

\section{3. Roteiro minuto a minuto}\label{roteiro-minuto-a-minuto}

\begin{longtable}[]{@{}
  >{\raggedright\arraybackslash}p{(\linewidth - 6\tabcolsep) * \real{0.0882}}
  >{\raggedright\arraybackslash}p{(\linewidth - 6\tabcolsep) * \real{0.0882}}
  >{\raggedright\arraybackslash}p{(\linewidth - 6\tabcolsep) * \real{0.1176}}
  >{\raggedright\arraybackslash}p{(\linewidth - 6\tabcolsep) * \real{0.7059}}@{}}
\toprule\noalign{}
\begin{minipage}[b]{\linewidth}\raggedright
Horário
\end{minipage} & \begin{minipage}[b]{\linewidth}\raggedright
Slides
\end{minipage} & \begin{minipage}[b]{\linewidth}\raggedright
Notebook
\end{minipage} & \begin{minipage}[b]{\linewidth}\raggedright
Atividade
\end{minipage} \\
\midrule\noalign{}
\endhead
\bottomrule\noalign{}
\endlastfoot
00:00 & 1--3 & & Apresentação, combinados e demonstração: jogar três
lances contra o modelo de referência \\
00:05 & 4 & Seção 0 & Acesso ao Colab: abrir o link, salvar cópia,
ativar a GPU, executar a seção 0. Aguardar a turma inteira \\
00:15 & 5--12 & & Motivação (RNNs, Tabela 1 do artigo) e atenção:
dicionário, Q/K/V, a equação, √d\_k, máscara causal \\
00:25 & 13 & Parte 1 & Exercícios 1 a 3 em NumPy. Aviso aos 00:40 \\
00:45 & 14--19 & & Multi-head, posição, bloco, arquitetura original,
famílias e o modelo completo \\
00:55 & 20 & Parte 2 & Exercícios 4 a 6 em PyTorch e treino da
ordenação \\
01:10 & 21 & 2.6 & Discutir o mapa de atenção da ordenação com a
turma \\
01:13 & 22 & 3.1, 3.2 & Todos iniciam o treino do modelo de xadrez \\
01:15 & 22 & & Intervalo de 5 minutos (o treino roda durante a pausa) \\
01:20 & 23--26 & & Xadrez como linguagem, objetivo de treino, modelos de
mundo, temperatura \\
01:30 & 27 & 3.3--3.6 & Gerar partidas, variar a temperatura, comparar
com o modelo de referência, jogar \\
01:50 & 28--30 & & Daqui até um LLM, referências e perguntas \\
02:00 & & & Encerramento \\
\end{longtable}

\textbf{Pontos de controle.} Em três momentos o ministrante deve
verificar o estado da sala antes de avançar: ao fim da seção 0 (todos
com GPU), aos 00:45 (todos com a Parte 1 concluída ou recuperada) e aos
01:13 (todos com o treino de xadrez em execução). Nesses momentos, peça
que quem ainda não chegou ao ponto execute as células de recuperação.

\textbf{Margem.} O roteiro não tem folga explícita. Se houver atraso,
reduza primeiro o slide 18 (famílias) e o slide 15 (posição), que podem
ser tratados em um minuto cada, e depois a seção 3.6 do notebook, que é
complementar.

\section{4. Condução da turma}\label{conduuxe7uxe3o-da-turma}

\textbf{Divisão da sala.} Com 40 ou mais participantes, divida o
laboratório em zonas, uma por monitor. O ministrante fica com a frente
da sala e circula apenas durante a Parte 2, que é a mais difícil.

\textbf{Sinal de ajuda.} Combine um sinal visível e silencioso, por
exemplo um papel dobrado sobre o monitor. Isso evita que o ministrante
seja interrompido durante a exposição e permite que os monitores
priorizem.

\textbf{Heterogeneidade.} Para estudantes dos primeiros semestres, a
orientação é explícita: o objetivo é entender o que cada exercício faz,
não necessariamente escrevê-lo sozinho. Após três minutos parados em um
exercício, eles devem usar a célula de recuperação e seguir. Incentive
duplas mistas (um estudante mais avançado e um iniciante), que funcionam
bem nesse formato.

\textbf{Estudantes adiantados.} Quem terminar antes pode fazer os
exercícios para casa do final do notebook já durante o workshop: trocar
o embedding de posição pelo senoidal ou aumentar o modelo.

\section{5. Erros comuns (para os
monitores)}\label{erros-comuns-para-os-monitores}

\begin{longtable}[]{@{}
  >{\raggedright\arraybackslash}p{(\linewidth - 4\tabcolsep) * \real{0.1148}}
  >{\raggedright\arraybackslash}p{(\linewidth - 4\tabcolsep) * \real{0.4098}}
  >{\raggedright\arraybackslash}p{(\linewidth - 4\tabcolsep) * \real{0.4754}}@{}}
\toprule\noalign{}
\begin{minipage}[b]{\linewidth}\raggedright
Exercício
\end{minipage} & \begin{minipage}[b]{\linewidth}\raggedright
Sintoma
\end{minipage} & \begin{minipage}[b]{\linewidth}\raggedright
Causa provável
\end{minipage} \\
\midrule\noalign{}
\endhead
\bottomrule\noalign{}
\endlastfoot
1 & ``a soma deve ser feita ao longo do último eixo'' & Falta
\texttt{keepdims=True} ou \texttt{axis=-1} \\
1 & NaN com entradas grandes & Não subtraiu o máximo antes de
\texttt{np.exp} \\
2 & Erro de shape em \texttt{Q\ @\ K} & Falta transpor \texttt{K} \\
2 & ``não foram divididos por √d\_k'' & Esqueceu a escala \\
3 & Pesos errados com máscara & Máscara com \texttt{np.triu} em vez de
\texttt{np.tril} \\
4 & Shape \texttt{(T,\ T,\ B)} ou similar & Usou \texttt{k.T}, que
inverte todas as dimensões em tensores 3D \\
4 & Erro de tamanho na máscara & Usou \texttt{self.tril} sem recortar
\texttt{{[}:T,\ :T{]}} \\
5 & Shape \texttt{(B,\ T,\ head\_size)} & Concatenou no eixo errado;
deve ser \texttt{dim=-1} \\
6 & ``o bloco devolveu a entrada sem alterações'' & As duas linhas não
foram escritas \\
6 & ``faltam as conexões residuais'' & Escreveu
\texttt{x\ =\ self.attn(...)} em vez de
\texttt{x\ =\ x\ +\ self.attn(...)} \\
Geral & \texttt{NameError} em células posteriores & A pessoa pulou uma
célula; executar tudo até o ponto atual (\emph{Ambiente de execução
\textgreater{} Executar anteriores}) \\
\end{longtable}

\section{6. Riscos e contingências}\label{riscos-e-continguxeancias}

\begin{longtable}[]{@{}
  >{\raggedright\arraybackslash}p{(\linewidth - 2\tabcolsep) * \real{0.2969}}
  >{\raggedright\arraybackslash}p{(\linewidth - 2\tabcolsep) * \real{0.7031}}@{}}
\toprule\noalign{}
\begin{minipage}[b]{\linewidth}\raggedright
Risco
\end{minipage} & \begin{minipage}[b]{\linewidth}\raggedright
Contingência
\end{minipage} \\
\midrule\noalign{}
\endhead
\bottomrule\noalign{}
\endlastfoot
Colab sem GPU disponível para parte da turma & Partes 1 e 2 funcionam em
CPU. Na Parte 3, a pessoa interrompe o treino e usa o modelo de
referência (seção 3.4), ou trabalha em dupla \\
Internet lenta no laboratório & Distribuir o repositório em pen drive
como plano B; o notebook funciona após \emph{Arquivos \textgreater{}
Upload} do zip e descompactação \\
Colab fora do ar & Projetar o gabarito já executado (salvo no ensaio) e
conduzir a Parte 3 como demonstração \\
Tabuleiro interativo não responde (falha na comunicação com o Colab) &
Em uma nova célula:
\texttt{from\ tabuleiro\ import\ jogar\_texto;\ jogar\_texto(ref,\ tok\_ref,\ device=device)},
versão com lances digitados \\
Modelo de referência ausente & O notebook usa automaticamente o modelo
treinado pela turma; a comparação da seção 3.4 fica prejudicada, mas
nada quebra \\
Atraso acumulado & Cortar conforme a seção 3 (slides 18 e 15, depois a
seção 3.6) \\
Turma muito heterogênea & Reforçar o uso das células de recuperação; os
monitores priorizam quem está parado na Parte 2 \\
\end{longtable}

\section{7. Estrutura do repositório}\label{estrutura-do-reposituxf3rio}

\begin{verbatim}
transformer-do-zero/
├── workshop_aluno.ipynb        notebook da turma (gerado)
├── workshop_gabarito.ipynb     soluções (gerado)
├── pretreino_referencia.ipynb  treino do modelo de referência (gerado)
├── mini_gpt.py                 implementação de referência
├── checagens.py                verificações dos exercícios
├── data/xadrez.txt.gz          conjunto de dados
├── modelos/                    xadrez_referencia.pt (após o pré-treino)
├── scripts/
│   ├── preparar_dados.py       PGN do Lichess -> data/xadrez.txt
│   ├── pretreinar.py           treino do modelo de referência
│   └── gerar_notebooks.py      fonte única dos notebooks
├── slides/
│   ├── attention_workshop.pptx
│   ├── gerar_slides.js         fonte da apresentação (pptxgenjs)
│   └── assets/                 equações e gráficos
└── docs/                       plano e guia de estudos (.md e .tex)
\end{verbatim}

Os notebooks e os slides são gerados a partir de scripts. Qualquer
alteração deve ser feita na fonte (\texttt{scripts/gerar\_notebooks.py}
ou \texttt{slides/gerar\_slides.js}) e o artefato regenerado. Editar o
\texttt{.ipynb} diretamente funciona, mas a alteração se perde na
próxima geração.

\textbf{Testes.} O comando abaixo executa o gabarito inteiro em modo
rápido (treinos de poucos passos, em CPU) e falha se alguma célula
levantar erro:

\begin{verbatim}
WORKSHOP_RAPIDO=1 .venv/bin/jupyter nbconvert --to notebook --execute \
    workshop_gabarito.ipynb --output /tmp/teste.ipynb
\end{verbatim}

\end{document}
